% =========================================================
\section{Post-contact degeneracy and fully clipped period-two dynamics}
\label{sec:post-contact}
% =========================================================

Sections~\ref{sec:switching-geometry} and
\ref{sec:separation-law} show that the smooth period-two branch
created by the supercritical flip reaches the switching hypersurface
simultaneously at both of its points.

The radial clipping structure imposes a further restriction:
a genuine period-two orbit cannot have one point in the smooth region
and the other in the clipped region. Hence any period-two orbit that
persists beyond the first switching contact must enter the clipped
regime at both points.

The post-contact results obtained below are not consequences of the
generic border-collision theory cited in
Section~\ref{sec:introduction}.  They rely specifically on the
radial structure of the clipping operator.

The purpose of this section is to characterize the universal
structure of such fully clipped period-two orbits.  The main result
is that every fully clipped period-two orbit is necessarily
nonhyperbolic: its two-step monodromy matrix has an exact multiplier
equal to one.

This neutral multiplier is not a consequence of the local flip
normal form.  It arises directly from the radial normalization
inherent in gradient clipping.

% ---------------------------------------------------------
\subsection{Fully clipped period-two orbits}
\label{subsec:fully-clipped-two-cycles}
% ---------------------------------------------------------

We first record the exact algebraic structure of a fully clipped
period-two orbit.

\begin{proposition}[Exact geometry of a fully clipped two-cycle]
\label{prop:fully-clipped-geometry}
Let
\[
\mathcal Y_2
=
\{y_-,y_+\}
\]
be a genuine period-two orbit of
\[
F_{\eta,c}(x)
=
x-\eta
\mathcal C_c^\diamond(\nabla V(x)),
\]
and suppose
\[
y_-,y_+\in\Omega_c^+.
\]
Then
\begin{equation}
\boxed{
\|y_+-y_-\|_\diamond
=
\eta c.
}
\label{eq:fully-clipped-step-length}
\end{equation}

Moreover, there exist
\[
u\in\mathbb R^d,
\qquad
\|u\|_\diamond=1,
\]
and numbers
\[
r_+,r_->c
\]
such that
\begin{align}
y_+-y_-
&=
\eta c\,u,
\label{eq:fully-clipped-displacement}
\\
\nabla V(y_+)
&=
r_+u,
\label{eq:fully-clipped-gradient-plus}
\\
\nabla V(y_-)
&=
-r_-u.
\label{eq:fully-clipped-gradient-minus}
\end{align}
\end{proposition}

\begin{proof}
The period-two relations are
\begin{align}
y_-
&=
y_+
-
\eta
\mathcal C_c^\diamond(\nabla V(y_+)),
\label{eq:fully-cycle-plus}
\\
y_+
&=
y_-
-
\eta
\mathcal C_c^\diamond(\nabla V(y_-)).
\label{eq:fully-cycle-minus}
\end{align}
Therefore
\begin{align}
\mathcal C_c^\diamond(\nabla V(y_+))
&=
\frac{y_+-y_-}{\eta},
\label{eq:fully-output-plus}
\\
\mathcal C_c^\diamond(\nabla V(y_-))
&=
-\frac{y_+-y_-}{\eta}.
\label{eq:fully-output-minus}
\end{align}

Since both points lie in the strictly clipped region,
\[
\left\|
\mathcal C_c^\diamond(\nabla V(y_\pm))
\right\|_\diamond
=
c.
\]
Taking norms in \eqref{eq:fully-output-plus} gives
\[
\|y_+-y_-\|_\diamond
=
\eta c.
\]

Define
\[
u
=
\frac{y_+-y_-}{\eta c}.
\]
Then
\[
\|u\|_\diamond=1.
\]

Because radial clipping preserves direction,
\[
\frac{\nabla V(y_+)}
{\|\nabla V(y_+)\|_\diamond}
=
u
\]
and
\[
\frac{\nabla V(y_-)}
{\|\nabla V(y_-)\|_\diamond}
=
-u.
\]
Setting
\[
r_\pm
=
\|\nabla V(y_\pm)\|_\diamond
\]
gives the result.
\end{proof}

\begin{remark}[Amplitude saturation]
\label{rem:amplitude-saturation-main}
Before contact, the period-two amplitude is generated by the smooth
flip mechanism and satisfies
\[
\|x_+(\mu)-x_-(\mu)\|
=
O(\sqrt{\mu}).
\]
After both points enter the clipped regime, the step length is no
longer determined by the smooth normal form.  Instead, radial clipping
imposes the exact constraint
\[
\|y_+-y_-\|_\diamond
=
\eta c.
\]

Thus clipping replaces the smooth square-root amplitude law by an
exact saturated step-length law.
\end{remark}

% ---------------------------------------------------------
\subsection{Exclusion of mixed period-two itineraries}
\label{subsec:no-mixed-two-cycle-main}
% ---------------------------------------------------------

We now state explicitly the structural restriction established in
Section~\ref{sec:switching-geometry}.

\begin{proposition}[No mixed smooth--clipped period-two orbit]
\label{prop:no-mixed-two-cycle-main}
A radial clipped gradient map cannot possess a genuine period-two
orbit
\[
\{y_-,y_+\}
\]
such that exactly one point lies in $\Omega_c^-$ and the other lies
in $\Omega_c^+$.

Consequently, every genuine period-two orbit is locally of one of the
following types:
\begin{enumerate}[label=\textnormal{(\roman*)}]
\item both points are smooth;
\item both points lie on the switching boundary;
\item both points are strictly clipped.
\end{enumerate}
\end{proposition}

\begin{proof}
Suppose, for contradiction, that
\[
y_-\in\Omega_c^-,
\qquad
y_+\in\Omega_c^+.
\]
Then the period-two equations imply
\[
\nabla V(y_-)
+
\mathcal C_c^\diamond(\nabla V(y_+))
=
0.
\]
Taking the clipping norm gives
\[
\|\nabla V(y_-)\|_\diamond
=
\left\|
\mathcal C_c^\diamond(\nabla V(y_+))
\right\|_\diamond
=
c,
\]
which contradicts
\[
\|\nabla V(y_-)\|_\diamond<c.
\]
The opposite mixed itinerary is ruled out identically.
\end{proof}

\begin{remark}
\label{rem:no-mixed-global}
Proposition~\ref{prop:no-mixed-two-cycle-main} is not a local
consequence of the flip bifurcation.  It follows from the algebraic
structure of radial clipping and therefore applies to arbitrary
period-two orbits, independently of their amplitude or origin.
\end{remark}

% ---------------------------------------------------------
\subsection{Differential structure in the clipped region}
\label{subsec:clipped-differential-main}
% ---------------------------------------------------------
We now examine the differential structure of the radial clipping
operator in the strictly clipped region. This structure will be
the key ingredient in the spectral analysis of fully clipped
period-two orbits.

Let
\[
\rho(g):=\|g\|_\diamond .
\]
For any nonzero vector $g\in\mathbb{R}^d$, write
\[
u:=\frac{g}{\rho(g)},
\qquad
\rho(u)=1,
\]
and define the covector
\[
\phi_u:=D\rho(u).
\]

Since $\rho$ is positively homogeneous of degree one,
Euler's identity gives
\begin{equation}
\phi_u[u]
=
D\rho(u)[u]
=
\rho(u)
=
1.
\label{eq:main-euler-rho}
\end{equation}
Moreover, the derivative of a positively homogeneous function
of degree one is homogeneous of degree zero. Hence, since
$g=\rho(g)u$ with $\rho(g)>0$,
\begin{equation}
D\rho(g)=D\rho(u)=\phi_u.
\label{eq:main-rho-derivative-homogeneity}
\end{equation}

In the strictly clipped region
\[
\rho(g)>c,
\]
the clipping operator takes the form
\[
C_c^\diamond(g)
=
c\frac{g}{\rho(g)}.
\]
For any $h\in\mathbb{R}^d$, differentiation yields
\[
DC_c^\diamond(g)[h]
=
\frac{c}{\rho(g)}h
-
\frac{c}{\rho(g)^2}
g\,D\rho(g)[h].
\]
Using
\[
g=\rho(g)u
\]
and \eqref{eq:main-rho-derivative-homogeneity}, we obtain
\[
DC_c^\diamond(g)[h]
=
\frac{c}{\rho(g)}
\left(
h-u\,\phi_u[h]
\right).
\]
Equivalently,
\begin{equation}
DC_c^\diamond(g)
=
\frac{c}{\rho(g)}
\left(
I-u\otimes\phi_u
\right).
\label{eq:main-clipping-derivative}
\end{equation}

The operator
\[
P_u:=I-u\otimes\phi_u
\]
has a simple geometric interpretation. By
\eqref{eq:main-euler-rho},
\[
P_u u
=
u-u\phi_u[u]
=
0,
\]
and
\[
\phi_u P_u
=
\phi_u-\phi_u[u]\phi_u
=
0.
\]
Thus $P_u$ projects onto
\[
T_u\{\rho=1\}
=
\ker\phi_u,
\]
the tangent space of the unit $\rho$-sphere at $u$, along the
radial direction spanned by $u$.

In particular, multiplying
\eqref{eq:main-clipping-derivative} from the left by $\phi_u$
gives
\begin{equation}
\phi_uDC_c^\diamond(g)
=
\frac{c}{\rho(g)}
\left(
\phi_u-\phi_u[u]\phi_u
\right)
=
0.
\label{eq:main-annihilation}
\end{equation}

Equation \eqref{eq:main-annihilation} expresses the differential
degeneracy created by radial clipping: infinitesimal variations
in the gradient-magnitude direction are annihilated, whereas only
variations tangent to the clipping sphere contribute to the
differential of the clipped update.

For every $x\in\Omega_c^+$, the chain rule therefore gives
\begin{equation}
DF_{\eta,c}(x)
=
I
-
\eta
DC_c^\diamond(\nabla V(x))
D^2V(x).
\label{eq:main-clipped-Jacobian}
\end{equation}

\begin{proposition}[Pointwise neutral covector]
\label{prop:pointwise-neutral-covector}
Let $x\in\Omega_c^+$ and set
\[
u
=
\frac{\nabla V(x)}
{\|\nabla V(x)\|_\diamond}.
\]
Then
\begin{equation}
\phi_uDF_{\eta,c}(x)=\phi_u.
\label{eq:main-pointwise-neutral}
\end{equation}
In particular,
\[
1\in\sigma\bigl(DF_{\eta,c}(x)\bigr).
\]
\end{proposition}

\begin{proof}
Using \eqref{eq:main-clipped-Jacobian},
\[
\phi_uDF_{\eta,c}(x)
=
\phi_u
-
\eta
\phi_u
DC_c^\diamond(\nabla V(x))
D^2V(x).
\]
By \eqref{eq:main-annihilation},
\[
\phi_u
DC_c^\diamond(\nabla V(x))
=
0.
\]
Hence the second term vanishes and
\[
\phi_uDF_{\eta,c}(x)=\phi_u.
\]
Thus $\phi_u$ is a left eigenvector of $DF_{\eta,c}(x)$
associated with the eigenvalue $1$. Therefore
\[
1\in\sigma\bigl(DF_{\eta,c}(x)\bigr).
\]
\end{proof}

\begin{remark}[Pointwise neutrality versus orbit neutrality]
\label{rem:pointwise-versus-orbit-neutrality}
Although every one-step Jacobian in the strictly clipped region
has a unit eigenvalue, the corresponding left eigenvector depends
on the local gradient direction. Hence, along a general clipped
orbit, unit eigenvalues of the individual Jacobians do not imply
that their product possesses a unit Floquet multiplier.

Period-two orbits are exceptional. As shown in
Proposition~7.1, the clipped period-two equations force the two
gradient directions to be exactly opposite. Moreover, because
$\rho$ is even,
\[
\rho(-u)=\rho(u),
\]
and therefore
\[
D\rho(-u)=-D\rho(u).
\]
Consequently,
\[
(-u)\otimes D\rho(-u)
=
u\otimes D\rho(u),
\]
so the tangent projection
\[
I-u\otimes\phi_u
\]
is unchanged under $u\mapsto -u$.

Thus the same neutral covector structure is preserved at both
points of a fully clipped period-two orbit. This observation is
the mechanism behind the universal unit-multiplier theorem proved
in the next subsection.
\end{remark}
                                                                                         


% ---------------------------------------------------------
\subsection{Universal unit multiplier for fully clipped two-cycles}
\label{subsec:universal-unit-main}
% ---------------------------------------------------------

We now obtain the principal post-contact result.

\begin{theorem}[Universal unit multiplier]

\label{thm:universal-unit-main}
Let
\[
\mathcal Y_2
=
\{y_-,y_+\}
\]
be any fully clipped period-two orbit of $F_{\eta,c}$.

Let
\[
M_c
:=
DF_{\eta,c}(y_-)
DF_{\eta,c}(y_+)
\label{eq:main-monodromy}
\]
be the corresponding two-step monodromy matrix.

Then
\begin{equation}
\boxed{
1\in\sigma(M_c).
}
\label{eq:main-unit-multiplier}
\end{equation}

More precisely, if $u$ is the common clipping direction defined by
Proposition~\ref{prop:fully-clipped-geometry}, then
\begin{equation}
\varphi_u
DF_{\eta,c}(y_+)
=
\varphi_u
\label{eq:common-left-plus}
\end{equation}
and
\begin{equation}
\varphi_u
DF_{\eta,c}(y_-)
=
\varphi_u.
\label{eq:common-left-minus}
\end{equation}
Hence
\begin{equation}
\varphi_uM_c
=
\varphi_u.
\label{eq:common-left-monodromy}
\end{equation}
\end{theorem}


\begin{proof}
By Proposition~\ref{prop:fully-clipped-geometry}, there exist
$r_+,r_->c$ and a vector $u$ with $\|u\|_\diamond=1$ such that
\[
\nabla V(y_+)=r_+u,
\qquad
\nabla V(y_-)=-r_-u.
\]

Since every norm is even,
\[
\rho(-z)=\rho(z),
\]
and therefore
\[
D\rho(-u)=-D\rho(u).
\]
Thus
\[
\phi_{-u}=-\phi_u.
\]

At $y_+$, the normalized gradient direction is $u$.
Hence Proposition~\ref{prop:pointwise-neutral-covector} gives
\[
\phi_u DF_{\eta,c}(y_+)=\phi_u.
\]

At $y_-$, the normalized gradient direction is $-u$.
Applying Proposition~\ref{prop:pointwise-neutral-covector}
with $-u$ gives
\[
\phi_{-u}DF_{\eta,c}(y_-)=\phi_{-u}.
\]
Since $\phi_{-u}=-\phi_u$, this is equivalent to
\[
\phi_uDF_{\eta,c}(y_-)=\phi_u.
\]

Consequently,
\[
\phi_uM_c
=
\phi_uDF_{\eta,c}(y_-)
DF_{\eta,c}(y_+)
=
\phi_u.
\]
Thus $\phi_u$ is a nonzero left eigenvector of $M_c$
associated with the eigenvalue $1$. Therefore
\[
1\in\sigma(M_c).
\]
\end{proof}

\begin{corollary}[Fully clipped period-two cycles are nonhyperbolic]
\label{cor:fully-clipped-nonhyperbolic}
Every fully clipped period-two orbit of a radial clipped gradient map
is nonhyperbolic.
\end{corollary}

\begin{proof}
Theorem~\ref{thm:universal-unit-main} gives a multiplier equal to
$1$, which lies on the unit circle.
\end{proof}

\begin{corollary}[No linear exponential attraction]
\label{cor:no-linear-exponential-attraction}
A fully clipped period-two orbit cannot be linearly exponentially
asymptotically stable.
\end{corollary}

\begin{proof}
Linear exponential asymptotic stability requires every Floquet
multiplier of the period-two orbit to lie strictly inside the unit
disk.  Theorem~\ref{thm:universal-unit-main} provides an exact
multiplier equal to $1$.
\end{proof}

\begin{remark}[Nonhyperbolic does not mean unstable]
\label{rem:nonhyperbolic-not-unstable}
The presence of the unit multiplier does not by itself imply
nonlinear instability.

A fully clipped period-two orbit may be nonlinearly attracting,
repelling, neutral, or embedded in a family of periodic orbits.
These alternatives depend on higher-order dynamics in the neutral
direction and cannot be distinguished by the linearization alone.
\end{remark}

% ---------------------------------------------------------
\subsection{The spectral transition at first contact}
\label{subsec:spectral-transition}
% ---------------------------------------------------------

Before contact, the smooth period-two orbit is hyperbolically
attracting along its critical direction.

By Theorem~\ref{thm:period-two-stability},
\begin{equation}
\rho_{\mathrm{sm}}^c(\mu)
=
1
-
4\lambda_*\mu
+
O(\mu^{3/2}).
\label{eq:precontact-critical-multiplier}
\end{equation}

At the first switching contact,
\begin{equation}
\mu_{\mathrm{sc}}(c)
=
K_\diamond c^2
+
O(c^4).
\label{eq:contact-mu-main}
\end{equation}
Hence
\begin{equation}
\rho_{\mathrm{sm}}^c
\bigl(
\mu_{\mathrm{sc}}(c)
\bigr)
=
1
-
4\lambda_*K_\diamond c^2
+
O(c^3).
\label{eq:contact-smooth-multiplier}
\end{equation}

Since
\[
K_\diamond
=
\frac{
\Gamma_*
}{
3\lambda_*^4
\|v_*\|_\diamond^2
},
\]
this becomes
\begin{equation}
\rho_{\mathrm{sm}}^c
\bigl(
\mu_{\mathrm{sc}}(c)
\bigr)
=
1
-
\frac{
4\Gamma_*
}{
3\lambda_*^3
\|v_*\|_\diamond^2
}
c^2
+
O(c^3).
\label{eq:contact-smooth-multiplier-Gamma}
\end{equation}

Thus the smooth cycle reaches the switching boundary while still
weakly attracting:
\[
\rho_{\mathrm{sm}}^c<1.
\]

Any fully clipped period-two orbit beyond contact, however, satisfies
\[
\rho_{\mathrm{clip}}^c=1
\]
in at least one direction.

This yields the following conclusion.

\begin{corollary}[Loss of hyperbolicity under fully clipped continuation]
\label{cor:loss-hyperbolicity-clipping}

Suppose that a period-two orbit admits a continuation from the
first switching contact into the fully clipped regime.
Then every such fully clipped period-two continuation is
necessarily nonhyperbolic.

Immediately before contact, the smooth critical multiplier satisfies
\[
\rho_{\mathrm{sm}}^c
=
1-4\lambda_*K_\diamond c^2+O(c^3),
\]
and therefore
\[
1-\rho_{\mathrm{sm}}^c
=
4\lambda_*K_\diamond c^2+O(c^3).
\]
Thus the smooth period-two branch reaches the switching boundary
while still weakly attracting, whereas, by
Theorem~\ref{thm:universal-unit-main}, every fully clipped
period-two orbit carries an exact unit multiplier.
\end{corollary}

\begin{remark}
\label{rem:spectral-transition-not-small-jump}
The transition should not be interpreted as a regular perturbation
of the smooth two-cycle multiplier.

The derivative of radial clipping changes discontinuously across
the switching boundary.  On the clipped side, radial variations of
the gradient are annihilated, and this rank-deficient structure
produces an exact neutral direction.

Thus the post-contact spectrum is constrained by the geometry of
the clipping operator rather than obtained by smoothly perturbing
the pre-contact smooth monodromy.
\end{remark}

% ---------------------------------------------------------
\subsection{Degeneracy of the period-two continuation equation}
\label{subsec:continuation-degeneracy}
% ---------------------------------------------------------

A period-two point of the clipped map is a solution of
\begin{equation}
\mathcal P(y,\eta,c)
=
F_{\eta,c}^2(y)-y
=
0.
\label{eq:period-two-continuation-equation}
\end{equation}

If the two-cycle were hyperbolic, one would normally attempt to
continue it by applying the implicit function theorem to
\[
\mathcal P=0.
\]
The universal unit multiplier prevents this.

\begin{proposition}[Singularity of the post-contact continuation problem]
\label{prop:continuation-singularity}
Let $y$ belong to a fully clipped period-two orbit. Then
\begin{equation}
D_y\mathcal P(y,\eta,c)
=
DF_{\eta,c}^2(y)-I
\label{eq:period-two-equation-Jacobian}
\end{equation}
is singular.

Consequently, assumptions \textbf{(H1)--(H5)} do not yield a
standard implicit-function continuation of the period-two branch
through the fully clipped regime.
\end{proposition}

\begin{proof}
By Theorem~\ref{thm:universal-unit-main},
\[
1
\in
\sigma
\left(
DF_{\eta,c}^2(y)
\right).
\]
Hence
\[
0
\in
\sigma
\left(
DF_{\eta,c}^2(y)-I
\right),
\]
so the derivative of the period-two fixed-point equation is
singular.
\end{proof}

\begin{remark}[Why a universal post-contact branch theorem is absent]
\label{rem:no-universal-postcontact-theorem}
The smooth period-two branch before contact is governed by a
nondegenerate flip bifurcation and therefore admits a standard local
branch description.

After clipping becomes active, the relevant periodic-orbit equation
is intrinsically degenerate.  Additional nonlinear information is
required to determine whether the fully clipped periodic orbit:

\begin{itemize}
\item persists as an isolated nonhyperbolic two-cycle;
\item belongs to a family of two-cycles;
\item is nonlinearly attracting or repelling;
\item or terminates at the switching contact.
\end{itemize}

The solvable planar model of Section~\ref{sec:planar-model} realizes
the second possibility: the post-contact dynamics contain a
one-parameter family of fully clipped period-two cycles.
\end{remark}

% ---------------------------------------------------------
\subsection{Euclidean clipping}
\label{subsec:euclidean-postcontact-main}
% ---------------------------------------------------------

For the Euclidean norm,
\[
\|\cdot\|_\diamond
=
\|\cdot\|_2,
\]
the geometry becomes especially transparent.

If
\[
u
=
\frac{\nabla V(x)}
{\|\nabla V(x)\|_2},
\]
then in the clipped region
\begin{equation}
D\mathcal C_c(\nabla V(x))
=
\frac{c}{\|\nabla V(x)\|_2}
\left(
I-uu^\top
\right).
\label{eq:euclidean-clipping-Jacobian-main}
\end{equation}

Hence
\begin{equation}
DF_{\eta,c}(x)
=
I
-
\frac{\eta c}{
\|\nabla V(x)\|_2
}
\left(
I-uu^\top
\right)
D^2V(x).
\label{eq:euclidean-clipped-map-Jacobian}
\end{equation}

The radial covector satisfies
\begin{equation}
u^\top DF_{\eta,c}(x)
=
u^\top.
\label{eq:euclidean-radial-left-eigenvector}
\end{equation}

For a fully clipped period-two orbit with gradient directions
$\pm u$,
\begin{equation}
u^\top
DF_{\eta,c}(y_-)
DF_{\eta,c}(y_+)
=
u^\top.
\label{eq:euclidean-two-step-neutral}
\end{equation}

Thus the unit multiplier has a direct geometric interpretation:
Euclidean clipping removes the derivative of the update in the
radial gradient-magnitude direction.

% ---------------------------------------------------------
\subsection{Summary of the post-contact mechanism}
\label{subsec:postcontact-summary-main}
% ---------------------------------------------------------

The local transition established in this paper can therefore be
summarized as follows.

For
\[
0<
\eta-\eta_f
<
\eta_{\mathrm{sc}}(c)-\eta_f,
\]
the period-two orbit lies entirely in the smooth region and is
governed by the supercritical smooth flip.

At
\[
\eta
=
\eta_{\mathrm{sc}}(c),
\]
both points of the orbit reach the switching boundary simultaneously.

A mixed smooth--clipped two-cycle cannot occur.

Any genuine period-two continuation into the clipped region must
satisfy
\[
\|y_+-y_-\|_\diamond
=
\eta c
\]
and necessarily possesses an exact unit multiplier:
\[
1
\in
\sigma(M_c).
\]

Hence the first clipping contact marks a transition
\begin{equation}
\begin{aligned}
\text{smooth attracting two-cycle}
&\\[-1mm]
\Downarrow
&\\[-1mm]
\text{simultaneous clipping contact}
&\\[-1mm]
\Downarrow
&\\[-1mm]
\text{nonhyperbolic clipped regime}.
\end{aligned}
\end{equation}
This degeneracy is a structural consequence of radial clipping and
does not depend on the specific solvable model considered later.